\lyseisE{Επανάληψη κεφαλαίου 5}
\solr{5}{1} $60\%$, $87{,}5\%$, $45\%$.
\solr{5}{2} $\frac{35}{100}=\frac{7}{20}$ και $\frac{12{,}5}{100}=\frac{1}{8}$.
\solr{5}{3} $36$ και $21$.
\solr{5}{4} $\frac{18}{72}=\frac{1}{4}=25\%$.
\solr{5}{5} $0{,}7\cdot140=98$ €.
\solr{5}{6} $1{,}12\cdot250=280$ €.
\solr{5}{7} $\frac{20}{80}=25\%$ αύξηση.
\solr{5}{8} $0{,}8x=64$, $x=80$ €.
\solr{5}{9} $1{,}24x=186$, $x=150$ €.
\solr{5}{10} Ένα έτος $60$ €, εννέα μήνες $0{,}75\cdot60=45$ €.
\solr{5}{11} Ο τέταρτος πήρε $10\%$. Ψήφοι: $180$, $120$, $60$ και $40$.
\solr{5}{12} $0{,}9\cdot0{,}9\cdot200=162$ €. Συνολική έκπτωση $38$ €, δηλαδή $19\%$ (όχι $20\%$).
\solr{5}{13} α) Σωστό. β) Σωστό. γ) Σωστό. δ) Λάθος: $10\permil=1\%$. ε) Σωστό ($10$ και $10$). στ) Λάθος: $0{,}5\%=0{,}005$.
\solr{5}{14} α) Λάθος: $44\%$. β) Λάθος: $0{,}9375$ της αρχικής. γ) Λάθος: τα $\frac{24}{124}$. δ) Σωστό. ε) Σωστό.
